% !TEX root = main.tex
\section{객체와 원장 상태}

\subsection{회계 값과 정규형}
각 객체의 회계 값은 질량 \(q\), 재활용 크레딧 \(a\), 탄소 \(e\)로 둔
다. \(q\)와 \(a\)는 같은 질량 단위와 배율을 사용하고, \(e\)는 고정 배
율의 kgCO2e 절대량이다. 모든 값은 음이 아닌 정수이며
\(0≤a≤q\)를 만족한다. \(a\)는 배분된 회계 크레딧이고 제품의 실
제 재활용 함량을 직접 나타내는 측정값으로 해석하지 않는다.
\(e\)는 등록 시 수용한 초기 탄소와 이후 증분·배분을 반영한다.

필드와 중간 계산에는 배포 설정에서 정한 범위를 적용한다.
합계·곱·나머지에도 범위를 검사하여 유한체의 모듈러 연산이
음수나 정수 오버플로(overflow)를 유효한 회계식으로 바꾸지

\begin{equation}
    \begin{gathered}y=\left\lfloor\frac{xq}{Q}\right\rfloor,\quad Q>0,\\xq=Qy+r,\quad 0\le r<Q.\end{gathered}\tag{3}
\end{equation}
분모가 0인 배분은 허용하지 않는다.

해시와 커미트먼트(commitment)는 타입, 필드 경계와 목적
을 구분하는 정규 직렬화와 도메인 분리(domain separation)
를 사용한다. 이하 \(H_{tag}(·)\)는 문자열 zkDPP:tag:v2를 도메
인으로 사용하는 해시의 축약이며, tag에는 Note, Voucher,
Nullifier 등 해당 목적의 식별자를 대입한다. 짧은 수학 표
기가 실제 직렬화 규칙을 대신하지는 않는다. 개방 난수와 소유
권 비밀의 영역, 유한체 원소로의 해시 변환(hash-to-field)과
문자열 정규화는 배포 설정에서 고정한다.

\subsection{회계 증표와 이전 증표}
회계 증표의 개방값과 공개 커미트먼트는

\begin{equation}
    \begin{gathered}n=(d,\rho,q,a,e,o,\omega_N),\\cm=H_{\mathrm{Note}}(d,\rho,q,a,e,o,\omega_N)\end{gathered}\tag{4}
\end{equation}
이다. 여기서 \(d\)는 문서 해시, \(ρ\)는 자산 역할, \(o\)는 소유자, \(ω_{N}\) 은
개방 난수다. 역할은 일반용(ELIGIBLE)과 폐기물(WASTE) 중

% Source PDF page 6, column 1
하나이다. WASTE 회계 증표는 \(a=e=0\)이며 종료(Exit)로만
소비한다. 소유권 비밀 \(sk\)로 계산하는 소비 식별자는

\begin{equation}
    nf=H_{\mathrm{Nullifier}}(cm,sk)\tag{5}
\end{equation}
이다. \(cm\)은 생성 때 공개하고, 미래 \(nf\)는 암호화하여 기록한
다. 실제 소비 시에는 같은 \(nf\)를 공개 입력으로 제출한다. 미래
식별자를 기록하는 것 자체가 소비를 뜻하지는 않는다.

이전 증표는 미해결 이전을 표현하며 다음 필드를 가진다.

\begin{equation}
    \begin{gathered}v=(d,\rho,q,a,e,o_S,o_R,D,\omega_V),\\rv=H_{\mathrm{Voucher}}(d,\rho,q,a,e,o_S,o_R,D,\omega_V).\end{gathered}\tag{6}
\end{equation}
\(o_{S},o_{R}\)는 송·수신 소유자이고 \(D\)는 절대 회수 기한이다. 이전
이벤트는 ELIGIBLE 이전 증표만 생성한다. 두 해결 경로가 사
용하는 공통 소비 식별자는

\begin{equation}
    \begin{gathered}s_{\mathrm{res}}=H_{\mathrm{VoucherResolutionSecret}}(\omega_V),\\rvnf=H_{\mathrm{VoucherNullifier}}(rv,s_{\mathrm{res}})\end{gathered}\tag{7}
\end{equation}
이다. 인수와 회수(Recall)는 각자의 소유권 지식을 증명하되
동일한 \(rvnf\)를 사용하므로, 원장에 먼저 수락된 경로만 이전
증표를 해결한다.

\subsection{문서, 제품 주장과 타입 참조}
문서 결합 대상은

\begin{equation}
    \mathrm{DocumentInfo}=(\mathrm{ProductName},\mathrm{LotID},\mathrm{Unit})\tag{8}
\end{equation}
이고 \(d=\) HashDocumentInfo\((\)DocumentInfo\()\)이다. 세 필드
는 DPP의 지정 경로에서 추출하며 누락·잘못된 타입·모호한 인
코딩을 거부한다. 발급(Issue) 정책 \(p\)와 공개 논스(nonce) \(c\)에
대해

\begin{equation}
    h=H_{\mathrm{Issue}}(d,p,c),\quad\mathrm{DPPClaim}=(p,h,c)\tag{9}
\end{equation}
를 정의한다. 발급은 회계 증표를 소비하여 공개 종단 객체 제
품 주장을 만든다. 제품 주장은 이후 입력으로 소비되지 않으며
소유자, 비공개 개방값, 포함 관계 트리, 미래 소비 식별자와 상
태가 없다.

객체는 타입을 가진 참조 \(r\)로 식별한다. 회계 증표에는
NoteRef\((cm)\), 이전 증표에는 VoucherRef\((rv)\), 제품 주장에
는 ClaimRef\((h)\)를 사용한다. 원장 조회 키는 ObjectRef 도메
인 안에 타입과 원시 식별자를 함께 결합하여 계산한다. 같은 숫
자를 갖더라도 회계 증표와 이전 증표를 같은 객체로 합치지 않
는다. 상태 조회에 쓰는 회계 증표의 \(nf\)와 이전 증표의 \(rvnf\)도
타입별로 구분한다.

\subsection{포함 관계, 소비 기록과 상태}
회계 증표와 이전 증표 커미트먼트는 서로 분리된 깊이 32의 추
가 전용 머클 트리(Merkle tree)에 등록한다. 노드, 리프(leaf)
수, 현재 루트(root)와 수락된 과거 루트를 원장에 보관하여 공
개 기록으로 포함 증명 경로를 재구성할 수 있도록 한다. 과거
루트의 포함 관계는 허용하되, 소비 승인에는 현재 소비 기록과
상태를 확인한다. 트리 용량을 넘는 추가는 전이 전체를 거부한
다.

% Source PDF page 6, column 2
수 락 된 이 벤 트 식 별 자 aid는 1 부 터 증 가 한 다. 타 입 별
spentIn\([ν]\)는 아직 소비되지 않았으면 0이고, 소비되었으면
해당 이벤트의 aid이다. 0에서 실제 소비 이벤트로 한 번 기록
한 값을 덮어쓰지 않는다. 공개 출력 참조에서 생성 이벤트로
가는 producerOf와 각 이벤트의 감사 기록도 수락 후 변경하지
않는다. 제품 주장은 별도의 등록 여부를 기록한다.

회계 증표와 이전 증표의 상태는 타입별 미래 소비 식별자
를 키로 하며 기본값은 활성(Active)이다. 나머지 상태는 동결
(Frozen)과 사용 금지(Revoked)이다. 허용되는 상태 변경은

\begin{equation}
    \begin{gathered}\mathrm{Active}\longrightarrow\mathrm{Frozen},\\\mathrm{Frozen}\longrightarrow\{\mathrm{Active},\mathrm{Revoked}\}\end{gathered}\tag{10}
\end{equation}
이다. Revoked는 종단이며 Active에서 바로 이동할 수 없다.
상태 관리자만 미소비 식별자에 대해 허용된 변경을 제출한다.
종료를 포함한 모든 소비는 현재 상태가 Active인 경우에만 허
용한다. 별도의 상태 트리나 비공개 상태 경로는 사용하지 않는
다. 매핑의 기본값만으로 객체의 존재가 입증되지는 않으며, 올
바른 객체의 식별자를 사용하는 책임은 신뢰된 감사자에게 있
다.

\subsection{전이와 결합된 감사 기록}
각 수락 이벤트는 종류, 정책 식별자, 순서 있는 공개 출력 참조,
부모 수, 실제 부모 목록의 암호문과 출력 미래 소비 식별자의
암호문을 포함한 감사 기록(AuditRecord)을 남긴다. 증명은 부
모 평문이 실제 포함 관계·소유권·회계 검사에 사용된 입력이
라는 사실과, 미래 식별자가 실제 출력 개방값에서 계산되었다
는 사실을 함께 검사해야 한다. 암호문은 지정된 감사 공개키
\(pk_{audit}\)와 문맥을 사용하며 배열의 타입·수·순서도 출력과 일치
해야 한다.

예를 들어 이전 이벤트의 출력 순서는 이전 증표, 잔여 회계
증표이고, 미래 식별자 암호문의 순서는 \(rvnf_{T}\) , \(nf_{C}\)이다. 제
품 주장에는 미래 소비 암호문이 없고 종료의 출력 관련 배열은
비어 있다. 초기 등록의 부모 목록은 빈 목록에 결합한다. 수락
시 배정되는 aid를 증명자가 미리 알아야 하는 암호화 문맥으로
요구하지 않으며, 실제 문맥은 공개 트랜잭션 자료로 재구성한
다.

\subsection{감사 결합의 관계 인터페이스}
앞의 일치 의무를 검토하기 위해 이벤트 종류 \(k\)의 공개 문장을
\(x=(u,ν,r,ct)\)로, 비공개 증인을 \(w=(z,sk,path,η)\)로
구분한다. \(u\)는 원장·배포 문맥, 이벤트 경로, 선택된 정책·권한,
포함 관계 루트와 필요한 기한 등 공개 메타데이터다. \(z\)는 실제
입력·출력 개방값, \(η\)는 암호화 난수다. 암호문을 계산하기 전의
공개 자료 \(u,ν,r\)와 슬롯의 목적·타입·위치로 \(ctx_{j}\)를 결정하며,
암호문 자체나 수락 후의 aid를 포함하지 않는다. 여기서 감사
공개키는 기록 암호화용 \(PK_{B}\)이며, 그 키 식별자도 \(ctx_{j}\)에 결
합한다. 감사 공개키와 그 식별자는 승인된 배포 설정에서 선택
하고, 증명자가 임의의 키로 대체할 수 없어야 한다. 배포에 키
가 고정되어 있으면 키 식별은 그 설정에서 결정할 수 있다. 데
이터 비밀키를 보호하는 \(W_{B}\)의 설정 검증은 외부 초기 설정의
책임이며, 이벤트 증인이 \(SK_{B}\)나 위원회의 비밀 조각을 포함하
도록 요구하지 않는다.

목적별 평문 함수 \(M_{j}(w)\)는 부모 슬롯이면 실제 입력의 타입
참조, 출력 슬롯이면 해당 출력의 타입과 미래 소비 식별자를 반

% Source PDF page 7, column 1
환한다. 등록된 관계는 적어도

\begin{equation}
    \begin{aligned}r_i^{\mathrm{out}}&=\operatorname{Ref}(z_i^{\mathrm{out}}),\\\nu_i^{\mathrm{in}}&=\operatorname{SpendId}(z_i^{\mathrm{in}},sk_i),\\ct_j&=\operatorname{AuditEnc}(pk_{\mathrm{audit}},M_j(w);\eta_j,\mathrm{ctx}_j)\end{aligned}\tag{11}
\end{equation}
를 개방·소유권·형상·회계 관계와 함께 검사해야 한다. 출력 슬
롯의 SpendId도 같은 출력 개방값과 적법한 소유권 비밀에서
계산한다. 복호화 정확성은 이 암호화 관계로 생성한 암호문에
대해 식 (2)가 \(M_{j}(w)\)를 반환한다는 계약이다. 임의의 잘못된
암호문을 검출할 능력과 구체 보안은 선택한 암호 방식의 별도
의무다.

소비 식별자는 서로 다른 객체의 구별뿐 아니라 동일 객
체의 효율적으로 생성·추출된 소비 증인 사이의 불변성도
요구한다. 즉 계산적 결합성의 나쁜 사건을 제외한 PPT 실
행에서 얻은, 같은 타입 참조를 여는 증인 \(w,w^{′}\)에 대해서는
SpendId\((w)=\) SpendId\((w^{′})\)이어야 하며, 그 값은 생성 시 암
호화한 식별자와 같아야 한다. 이를 위해 개방값의 결합성과 함
께 소유권 비밀의 정규 영역 및 OwnerKey의 증인 별칭을 배제
하는 성질이 필요하다. 예를 들어 OwnerKey\((s)=s\) mod \(5\)인
데 1과 6을 모두 유효 비밀로 허용하면, 같은 소유자라도 식 (5)의
두 식별자가 달라질 수 있다. 이 예는 구체 구현의 취약점이 아
니라 인터페이스의 불변성 가정을 생략할 수 없음을 보인다. 이
전 증표에도 개방 난수의 결합성과 두 해결 경로의 같은 식별자
계산을 요구한다. 이하 분석은 이 관계 계약을 만족하는 실행
에 대한 것이며 특정 곡선·회로가 이를 달성했다는 주장은 아니
다.
